\documentclass[preprint,review,12pt]{elsarticle}

%% Packages
\usepackage{lineno}
\usepackage{hyperref}
\usepackage{amsmath,amssymb}
\usepackage{booktabs}
\usepackage{graphicx}
\usepackage{multirow}
\usepackage{xcolor}
\usepackage{siunitx}
\usepackage{textcomp}

%% Journal name
\journal{Environmental Modelling \& Software}

\begin{document}

\begin{frontmatter}

\title{EVAonline: A web-based platform for global reference evapotranspiration estimation through adaptive Kalman filter fusion of multi-source climate data}

%% Authors
\author[unesp]{Angela Cunha Soares\corref{cor1}}
\ead{angela.soares@unesp.br}
%% TODO: Add additional authors as needed
%% \author[inst2]{Second Author}
%% \ead{second.author@email.com}

\cortext[cor1]{Corresponding author.}

%% Addresses
\address[unesp]{S\~ao Paulo State University (UNESP), Brazil}
%% \address[inst2]{Institution, Country}

%% Abstract
\begin{abstract}
Reference evapotranspiration (ET\textsubscript{0}) is a critical parameter for
irrigation scheduling, hydrological modelling, and climate impact assessment,
yet its accurate estimation remains challenging in regions with sparse
meteorological networks. This paper presents EVAonline, an open-source,
web-based platform that improves ET\textsubscript{0} estimation through a
two-stage data fusion strategy: (i)~region-adaptive weighted averaging of six
public climate APIs (NASA POWER, Open-Meteo Archive, Open-Meteo Forecast, MET
Norway, NWS Daily Forecast, and NWS Stations), followed by (ii)~an adaptive
one-dimensional Kalman filter constrained by 30-year regional monthly
climatology (1991--2020 normals and P01/P99 percentile bounds). The fused
meteorological variables serve as input to a complete FAO-56 Penman-Monteith
implementation. The platform was validated against the BR-DWGD gridded
reference dataset at 17 Brazilian sites (16 in the MATOPIBA agricultural
frontier plus Piracicaba/SP), yielding a Kling-Gupta Efficiency (KGE) of $0.814 \pm
0.053$, Nash-Sutcliffe Efficiency (NSE) of $0.676 \pm 0.085$, and Mean
Absolute Error (MAE) of \SI{0.423}{mm.d^{-1}}, with a mean bias (PBIAS) of
only $0.71 \pm 0.53\%$. The fusion reduced RMSE by 49\% relative to
NASA POWER (from 1.117 to 0.566~mm\,d$^{-1}$) and by 34\% relative to the
best individual source (Open-Meteo API).
EVAonline is deployed as 13 Docker-containerized microservices following a
hexagonal (Ports \& Adapters) architecture, requires no local installation or
API keys, and provides three operational modes (historical, recent, and
forecast), interactive dashboards, CSV/Excel data exports via email, and a
bilingual interface (English/Portuguese). The source code is available under
the AGPL-3.0 license at \url{https://github.com/angelacunhasoares/EVAONLINE},
and the complete validation dataset is published on Zenodo.
\end{abstract}

%%Graphical abstract
\begin{graphicalabstract}
%% TODO: Create a graphical abstract showing the pipeline:
%% Left: 6 climate API icons → Center: Two-stage fusion (weighted avg + Kalman)
%% → Right: FAO-56 PM → Dashboard screenshot + validation scatter plot
%% Recommended dimensions: 525 × 292 pixels minimum, landscape orientation
%\includegraphics[width=\textwidth]{graphical_abstract}
\end{graphicalabstract}

%%Research highlights
\begin{highlights}
\item EVAonline fuses six public climate APIs for global ET\textsubscript{0} estimation
\item Two-stage fusion: weighted averaging plus adaptive Kalman filter
\item Validated at 17 MATOPIBA sites achieving KGE = 0.814
\item Open-source Docker platform under AGPL-3.0 with web dashboard
\item No API keys or local installation required
\end{highlights}

%% Keywords
\begin{keyword}
Reference evapotranspiration \sep Data fusion \sep Adaptive Kalman filter \sep
FAO-56 Penman-Monteith \sep Open-source software \sep Web application \sep
Irrigation management
\end{keyword}

\end{frontmatter}

\linenumbers

%%
%% INTRODUCTION
%%
\section{Introduction}
\label{sec:introduction}

Reference evapotranspiration (ET\textsubscript{0}) is a fundamental parameter
in hydrology, agricultural water management, irrigation scheduling, and
climate impact assessment. The FAO-56 Penman-Monteith (PM) equation
\citep{Allen1998} is the internationally recommended standard for
ET\textsubscript{0} estimation, requiring complete meteorological datasets
including maximum and minimum air temperature, relative humidity, wind speed,
solar radiation, and atmospheric pressure. Despite its widespread adoption, the
application of the FAO-56 PM equation is frequently constrained by the
availability, quality, and spatial density of meteorological observations,
particularly in developing countries and remote agricultural regions.

A persistent challenge in ET\textsubscript{0} estimation is the reliability of
input meteorological data in regions with sparse weather station networks.
Gridded reanalysis and satellite-derived datasets---such as NASA POWER
\citep{Stackhouse2018}, ERA5 \citep{Hersbach2020}, and numerical weather
prediction (NWP) outputs---offer global coverage but exhibit source-specific
biases and uncertainties that vary by variable, region, and season. Using a
single data source introduces systematic errors that propagate through the
nonlinear PM equation, potentially leading to significant ET\textsubscript{0}
estimation errors \citep{Bai2010, Blankenau2020, Paredes2020}. Multiple
studies have demonstrated that no single gridded product consistently
outperforms others across all meteorological variables and climatic regions
\citep{Monteiro2022, Sheffield2006}, suggesting that data fusion approaches
may yield more robust estimates.

The MATOPIBA region---comprising the Brazilian states of Maranh\~ao (MA),
Tocantins (TO), Piau\'i (PI), and western Bahia (BA)---exemplifies this
challenge. MATOPIBA is one of the most dynamic agricultural frontiers in the
world, responsible for approximately 10\% of Brazilian grain production, with
cropland area expanding at rates exceeding 3\% per year
\citep{Miranda2014, Bolfe2016}. The region's tropical climate (K\"oppen
Aw/BSh), with a pronounced dry season from May to September and mean annual
temperatures of 25--28~\textdegree{}C, makes accurate
ET\textsubscript{0} estimation critical for irrigation planning and water
resource management. However, MATOPIBA has one of the lowest densities of
meteorological stations in Brazil \citep{Xavier2016}, creating a significant
gap between the precision required for irrigation management and the
observational data available.

Several software tools exist for ET\textsubscript{0} computation: CROPWAT~8.0
\citep{cropwat} and the ETo Calculator \citep{etocalculator} from FAO are
widely used desktop applications that implement the FAO-56 PM equation but
require manual data input and offer no multi-source fusion capabilities;
Ref-ET \citep{refet} supports multiple equation formulations but is a
commercial desktop product; PyETo \citep{pyeto} is an open-source Python
library that provides programmatic access but requires programming knowledge;
and OpenET \citep{openet} is a web-based platform focused on actual
evapotranspiration from satellite imagery over the contiguous United States.
None of these tools combines automatic multi-source data acquisition, optimal
data fusion with quantified uncertainty, web-based access requiring no
installation, and global coverage in a single, freely available platform.

This paper presents EVAonline, a web-based platform that addresses these
limitations through an integrated approach combining: (a)~automatic data
acquisition from six public climate APIs covering both historical reanalysis
and real-time forecast products; (b)~a two-stage data fusion strategy
consisting of region-adaptive weighted averaging followed by adaptive Kalman
filtering with climatological constraints; (c)~a complete FAO-56
Penman-Monteith implementation with water deficit analysis; and (d)~a
production-ready web interface accessible through any modern browser,
eliminating the need for local software installation, programming knowledge,
or manual data preprocessing. EVAonline is implemented as 13 Docker-containerized
microservices following a hexagonal architecture, ensuring modularity, testability,
and straightforward deployment.

The remainder of this paper is organized as follows:
Section~\ref{sec:methodology} describes the system architecture, data sources,
fusion methodology, and ET\textsubscript{0} calculation procedures;
Section~\ref{sec:results_discussion} presents the validation results against
the BR-DWGD reference dataset, compares EVAonline with existing tools, and
discusses strengths and limitations; and Section~\ref{sec:conclusions}
summarizes the main contributions and outlines future development directions.


%%
%% METHODOLOGY
%%
\section{Methodology}
\label{sec:methodology}

\subsection{System Architecture}
\label{subsec:architecture}

EVAonline is implemented as a distributed microservice architecture comprising
13 containerized services orchestrated through Docker Compose
(Fig.~\ref{fig:architecture}). The system follows a hexagonal (Ports \&
Adapters) architecture pattern \citep{Cockburn2005} that enforces strict
separation between domain logic, external service integration, and
infrastructure concerns.

The architecture is organized into three concentric layers:

\textbf{Domain core.} The innermost layer contains all business logic with zero
external I/O dependencies, including the FAO-56 PM calculation service
(\texttt{EToCalculationService}), the multi-source data fusion engine
(\texttt{ClimateFusion}), and the adaptive Kalman filter
(\texttt{AdaptiveKalmanFilter}). This layer operates exclusively on in-memory
data structures (Pandas DataFrames and NumPy arrays), ensuring testability and
portability.

\textbf{Adapter layer.} Six climate data adapters implement a standardized
interface for external API communication. Each adapter comprises a client class
(handling HTTP requests, retries, and response parsing) and a synchronization
adapter (transforming source-specific schemas into the standardized internal
representation). A \texttt{ClimateClientFactory} provides dependency injection,
enabling runtime selection of data sources based on geographic region and
operational mode.

\textbf{Infrastructure layer.} External services---Redis 7.x for caching and
message brokering, PostgreSQL 16 with PostGIS for spatial data persistence,
and Celery 5.x for distributed task processing---are accessed exclusively
through infrastructure adapters, ensuring the domain core remains independent
of specific technology choices.

The frontend is built with Dash~2.x \citep{plotly} and communicates with the
FastAPI backend \citep{Ramirez2018} exclusively through RESTful HTTP endpoints
and WebSocket connections. Nginx serves as a reverse proxy handling SSL/TLS
termination, rate limiting (30~requests/s general, 5~requests/s for
computations), and static asset serving. Application monitoring is provided by
a Prometheus + Grafana stack that collects request latency, error rates, and
calculation throughput metrics, while Celery task monitoring is provided
through Flower.

Deployment requires a single \texttt{docker compose up} command, which
initializes the complete platform including database creation, schema migration
via Alembic, seed data loading, and service health verification. The
multi-stage Docker build produces a production image of approximately 450~MB.

\begin{figure}[htbp]
\centering
%% TODO: Insert architecture diagram
%\includegraphics[width=0.95\textwidth]{fig_architecture}
\caption{System architecture of EVAonline showing the 13 Docker-containerized
services organized in a hexagonal (Ports \& Adapters) pattern. The domain core
(center) contains FAO-56 PM calculation and Kalman fusion logic with no
external dependencies. Six climate adapters (left) and infrastructure services
(right) communicate through well-defined interfaces.}
\label{fig:architecture}
\end{figure}


\subsection{Climate Data Sources and Variable Standardization}
\label{subsec:data_sources}

EVAonline integrates meteorological data from six public APIs provided by four
organizations (Table~\ref{tab:sources}). Each API is encapsulated as a Python
adapter class that enforces a standardized data retrieval interface with
configurable HTTP timeouts and exponential-backoff retry logic via the
\texttt{httpx} asynchronous HTTP client.

\begin{table}[H]
\caption{Climate data sources integrated in EVAonline.}
\label{tab:sources}
\resizebox{\textwidth}{!}{%
\begin{tabular}{lllc}
\hline
\textbf{Source}     & \textbf{Model}    & \textbf{Coverage}     & \textbf{Resolution}  \\ \hline
NASA POWER v2.0     & MERRA-2 + CERES   & Global, 1981--present & 0.5$^{\circ}$        \\
Open-Meteo Archive  & Best Match$^b$ (IFS/ERA5) & Global, 1940--present & 9--25\,km    \\
Open-Meteo Forecast & Best Match$^b$ (multi-NWP) & Global, 16 days      & 1--25\,km    \\
MET Norway (Yr)     & MEPS / ECMWF IFS  & Global, 10 days       & 2.5--9\,km$^a$       \\
NWS Forecast        & NWP (GFS/NAM)     & CONUS, 7 days         & 2.5\,km              \\
NWS Stations        & In-situ obs.      & CONUS ($\sim$1800 st.)& Station              \\ \hline
\end{tabular}%
}
\begin{flushleft}
\scriptsize
$^a$~2.5\,km in Nordic Europe (MEPS); $\sim$9\,km elsewhere (ECMWF IFS).\\
$^b$~Open-Meteo \texttt{best\_match} auto-selects the highest-resolution model per location.
\textit{Archive}: ECMWF IFS (9\,km, 2017+) or ERA5 (25\,km, pre-2017).
\textit{Forecast}: combines ICON, GFS, ECMWF IFS, AROME, and 10+ regional NWP models.\\
EVAonline enforces standardised limits: historical queries start from 1990-01-01;
forecast horizon is capped at today\,+\,5\,d (6 days total) across all sources.
\end{flushleft}
\end{table}

Source-specific variable names, units, and temporal aggregations differ across
APIs. Each adapter transforms the native schema into a standardized internal
representation using the seven meteorological variables required by the FAO-56
PM equation (Table~\ref{tab:variables}). After standardization, all variables
share a common naming convention, SI-compatible units, and daily temporal
resolution.

\begin{table}[htbp]
\centering
\caption{Variable standardization across data sources. Each source maps its
native variable names to a common internal schema.}
\label{tab:variables}
\small
\begin{tabular}{@{}lllll@{}}
\toprule
Internal variable & Unit & NASA POWER & Open-Meteo & MET Norway \\
\midrule
T\textsubscript{max} (\texttt{T2M\_MAX}) & \textdegree{}C & \texttt{T2M\_MAX} & \texttt{temperature\_2m\_max} & daily max \\
T\textsubscript{min} (\texttt{T2M\_MIN}) & \textdegree{}C & \texttt{T2M\_MIN} & \texttt{temperature\_2m\_min} & daily min \\
T\textsubscript{mean} (\texttt{T2M}) & \textdegree{}C & \texttt{T2M} & \texttt{temperature\_2m\_mean} & computed \\
RH (\texttt{RH2M}) & \% & \texttt{RH2M} & \texttt{relative\_humidity\_2m\_mean} & \texttt{relative\_humidity} \\
$u_z$ (\texttt{WS2M}) & m\,s\textsuperscript{$-1$} & \texttt{WS10M} ($z$=10\,m) & \texttt{wind\_speed\_10m\_max} & \texttt{wind\_speed} \\
$R_s$ (\texttt{ALLSKY\_SFC\_SW\_DWN}) & MJ\,m\textsuperscript{$-2$}\,d\textsuperscript{$-1$} & \texttt{ALLSKY\_SFC\_SW\_DWN} & \texttt{shortwave\_radiation\_sum} & N/A \\
Precipitation (\texttt{PRECTOTCORR}) & mm\,d\textsuperscript{$-1$} & \texttt{PRECTOTCORR} & \texttt{precipitation\_sum} & \texttt{precipitation\_amount} \\
\bottomrule
\end{tabular}
\end{table}

Site elevation, required for the barometric estimation of atmospheric pressure,
is obtained from the Open-Elevation API (SRTM 30\,m resolution). When the API
is unavailable, a latitude-based empirical function provides a fallback
estimate. Elevation values are cached in Redis with a 24-hour time-to-live
(TTL). Atmospheric pressure is estimated from elevation using the barometric
formula \citep{Allen1998}:

\begin{equation}
P = 101.3 \left(\frac{293 - 0.0065\,z}{293}\right)^{5.26}
\label{eq:pressure}
\end{equation}

\noindent where $z$ (m) is the site elevation and $P$ (kPa) is the atmospheric
pressure.

All wind speed measurements are adjusted to the standard reference height of
2~m using the FAO-56 logarithmic wind profile equation \citep{Allen1998}:

\begin{equation}
u_2 = u_z \cdot \frac{4.87}{\ln(67.8\,z - 5.42)}
\label{eq:wind}
\end{equation}

\noindent where $u_z$ (m\,s\textsuperscript{$-1$}) is the wind speed measured
at height $z$ (m). For NASA POWER and Open-Meteo ($z = 10$~m), the correction
factor is 0.748.


\subsection{Data Quality Control and Multi-Source Fusion}
\label{subsec:fusion}

The EVAonline data processing pipeline employs a two-stage fusion strategy:
region-adaptive weighted averaging of standardized multi-source data, followed
by adaptive Kalman filtering constrained by 30-year climatological priors
(Fig.~\ref{fig:fusion}). Prior to fusion, all data undergo a multi-layer
quality control procedure.

\begin{figure}[htbp]
\centering
%% TODO: Insert fusion pipeline flowchart
%\includegraphics[width=0.90\textwidth]{fig_fusion_pipeline}
\caption{Two-stage data fusion pipeline. Stage~1 performs quality control and
region-adaptive weighted averaging across available sources. Stage~2 applies an
adaptive Kalman filter to precipitation and ET\textsubscript{0}, using monthly
climatological priors (normals and P01/P99 percentile bounds) from the
1991--2020 reference period.}
\label{fig:fusion}
\end{figure}

\subsubsection*{Quality control}

Each source's data undergoes three sequential quality control steps before
entering the fusion stage:

\begin{enumerate}
\item \textbf{Physical range validation.} Measurements are checked against
  region-specific physically plausible bounds derived from WMO records and
  the Brazilian weather gridded dataset \citep{Xavier2016, Xavier2022}. For
  Brazilian locations, tighter bounds are applied (e.g., temperature
  $-30$~to~50~\textdegree{}C, precipitation 0--450~mm\,d\textsuperscript{$-1$},
  solar radiation 0--40~MJ\,m\textsuperscript{$-2$}\,d\textsuperscript{$-1$}).
  Global bounds are wider to accommodate extreme climates (e.g., temperature
  $-90$~to~60~\textdegree{}C). Values exceeding these ranges are set to NaN.

\item \textbf{Outlier detection.} An adaptive interquartile range (IQR) method
  identifies statistical outliers with a configurable IQR factor. A maximum of
  5\% of records may be flagged as outliers per variable to prevent excessive
  data removal. Solar radiation is additionally validated against
  extraterrestrial radiation ($R_a$) to ensure $0.03\,R_a \leq R_s \leq R_a$.

\item \textbf{Circuit-breaker mechanism.} A per-source quality score is
  computed as the product of completeness (fraction of non-null values) and
  validity (fraction within physical bounds) for each variable. Sources with an
  average quality score below 60\% are automatically excluded from the fusion
  for the current request, preventing degraded data from contaminating the
  fused output.
\end{enumerate}

\subsubsection*{Stage 1: Region-adaptive weighted averaging}

For each day $k$ and each meteorological variable $j$, the Stage~1 fused
estimate is computed as a normalized weighted average of all healthy sources
providing valid data:

\begin{equation}
\hat{z}_{k,j} = \frac{\sum_{i \in \mathcal{S}_{k,j}} w_{i,j} \cdot z_{k,i,j}}{\sum_{i \in \mathcal{S}_{k,j}} w_{i,j}}
\label{eq:weighted_avg}
\end{equation}

\noindent where $\mathcal{S}_{k,j}$ is the set of sources providing a valid
measurement for variable $j$ on day $k$, $z_{k,i,j}$ is the observation from
source $i$, and $w_{i,j}$ are source and variable-specific weights.
When only a single source provides data for a given variable on a given day,
that source's value is used directly ($w_{i,j} = 1$). Two weighting schemes
are employed depending on the operational mode.

\paragraph{Historical mode weights.}
For the two-source historical configuration (NASA POWER + Open-Meteo Archive),
variable-specific weights are assigned based on the comparative performance of
each reanalysis product against the BR-DWGD reference dataset across the 17
validation sites (Table~\ref{tab:hist_weights}). NASA POWER receives
higher weight for solar radiation ($w = 0.92$), leveraging its CERES satellite
observations, while Open-Meteo (ERA5) receives higher weight for wind speed
($w_{\text{OM}} = 0.80$) and relative humidity ($w_{\text{OM}} = 0.65$).
Precipitation is computed as the simple arithmetic mean of all available
sources, since the subsequent Kalman filter (Stage~2) provides the primary
quality control for this variable.

\begin{table}[htbp]
\centering
\caption{Variable-specific source weights for historical mode (NASA POWER +
Open-Meteo Archive). The NASA weight $w_{\text{NASA}}$ is shown; Open-Meteo
receives $1 - w_{\text{NASA}}$.}
\label{tab:hist_weights}
\small
\begin{tabular}{@{}lccc@{}}
\toprule
\textbf{Variable} & $w_{\text{NASA}}$ & $w_{\text{OM}}$ & \textbf{Rationale} \\
\midrule
$T_{\max}$    & 0.58 & 0.42 & NASA slight advantage in extremes \\
$T_{\min}$    & 0.52 & 0.48 & Near-equal performance \\
$T_{\text{mean}}$ & 0.60 & 0.40 & NASA MERRA-2 advantage \\
RH        & 0.35 & 0.65 & ERA5 humidity superior \\
$u_2$     & 0.20 & 0.80 & ERA5 wind more accurate \\
$R_s$     & 0.92 & 0.08 & CERES satellite radiation dominant \\
Precipitation & \multicolumn{2}{c}{simple mean} & Kalman applied after \\
\bottomrule
\end{tabular}
\end{table}

\paragraph{Forecast and recent mode weights.}
For forecast and recent periods with up to three real-time sources,
region-dependent uniform weights per source are applied
(Table~\ref{tab:weights}).

\begin{table}[htbp]
\centering
\caption{Region-dependent source weights for forecast and recent mode
weighted averaging. The geographic region is determined automatically from
the query coordinates.}
\label{tab:weights}
\begin{tabular}{@{}llccc@{}}
\toprule
Region & Detection criteria & \multicolumn{3}{c}{Source weights} \\
\cmidrule(l){3-5}
 & & NWS Forecast & Open-Meteo Forecast & MET Norway \\
\midrule
Americas (USA) & $24\leq\varphi\leq50$, $-125\leq\lambda\leq-66$ & 0.50 & 0.30 & 0.20 \\
Nordic Europe & $55\leq\varphi\leq72$, $-10\leq\lambda\leq40$ & --- & 0.20 & 0.80 \\
Global (default) & All other locations & --- & 0.70 & 0.30 \\
\bottomrule
\end{tabular}
\begin{flushleft}
\small Notes: NWS sources are available only within the continental United
States. $\varphi$: latitude (\textdegree{}N); $\lambda$: longitude
(\textdegree{}E).
\end{flushleft}
\end{table}

After weighted averaging, the fused values are clipped to physical bounds and
gaps of up to three consecutive days are filled by linear interpolation,
constrained within the variable-specific physical limits to prevent
extrapolation artifacts. Days with fewer than four valid variables are discarded.

\subsubsection*{Stage 2: Adaptive Kalman smoothing}

Following the weighted averaging stage, an adaptive one-dimensional scalar
Kalman filter \citep{Kalman1960} is applied in two sequential phases: first to
precipitation (immediately after Stage~1, before the FAO-56 calculation), and
then to ET\textsubscript{0} (after the FAO-56 Penman-Monteith calculation,
Section~\ref{subsec:eto}). Both applications use the same filter formulation
with variable-specific climatological priors, but differ in initialization and
reset strategy.

For each filtered variable $v \in \{\text{precipitation}, \text{ET}_0\}$ and
each calendar month $m \in \{1, \ldots, 12\}$, a separate filter instance is
initialized using climatological statistics from the 1991--2020 reference
period:

\begin{itemize}
\item Prior state estimate: $\hat{x}_0 = \mu_m$ (climatological monthly normal)
\item Prior error variance: $P_0 = \sigma_m^2$ (climatological monthly variance,
  enforced minimum of 0.5~mm\,d\textsuperscript{$-1$} for ET\textsubscript{0}
  and 5.0~mm\,d\textsuperscript{$-1$} for precipitation)
\item Process noise: $Q_m = 0.08\,\sigma_m^2$
\item Base measurement noise: $R_{\text{base}} = \sigma_{\text{meas}}^2 =
  0.55^2 = 0.3025$, where $\sigma_{\text{meas}} = 0.55$~mm\,d\textsuperscript{$-1$}
  is the assumed measurement standard deviation
\end{itemize}

The state equation follows a random walk model appropriate for day-to-day
meteorological persistence:

\begin{equation}
x_k = x_{k-1} + w_{k-1}, \quad w_{k-1} \sim \mathcal{N}(0, Q_m)
\label{eq:kalman_state}
\end{equation}

\noindent with the observation model:

\begin{equation}
z_k = x_k + v_k, \quad v_k \sim \mathcal{N}(0, R_k)
\label{eq:kalman_obs}
\end{equation}

\noindent where $z_k$ is the Stage~1 fused estimate for day $k$, and $R_k$ is
an adaptive measurement noise variance that increases when observations deviate
from climatologically expected ranges. Specifically:

\begin{equation}
R_k =
\begin{cases}
500\,R_{\text{base}} & \text{if } z_k < 0.8\,\text{P01}_m \text{ or } z_k > 1.25\,\text{P99}_m \\
50\,R_{\text{base}} & \text{if } z_k < \text{P01}_m \text{ or } z_k > \text{P99}_m \\
R_{\text{base}} & \text{otherwise}
\end{cases}
\label{eq:adaptive_R}
\end{equation}

\noindent where $\text{P01}_m$ and $\text{P99}_m$ are the 1st and 99th
percentile bounds of the local climatological distribution for calendar
month~$m$. When P01$_m$ or P99$_m$ are unavailable (e.g., for locations
outside the reference city network), they default to
$\mu_m \pm 3.5\,\sigma_m$, providing conservative bounds that still enable
anomaly detection. This three-tier scheme effectively downweights observations
that lie outside the climatologically plausible range, with extreme outliers
(beyond $0.8 \times \text{P01}$ or $1.25 \times \text{P99}$) receiving
500-fold noise inflation---retaining their partial information rather than
discarding them entirely.

The standard Kalman prediction and update equations are applied:

\begin{align}
\hat{x}_k^{-} &= \hat{x}_{k-1}^{+} \label{eq:kalman_predict_x}\\
P_k^{-} &= P_{k-1}^{+} + Q_m \label{eq:kalman_predict_P}\\
K_k &= \frac{P_k^{-}}{P_k^{-} + R_k} \label{eq:kalman_gain}\\
\hat{x}_k^{+} &= \hat{x}_k^{-} + K_k\,(z_k - \hat{x}_k^{-}) \label{eq:kalman_update_x}\\
P_k^{+} &= (1 - K_k)\,P_k^{-} \label{eq:kalman_update_P}
\end{align}

\noindent where $K_k$ is the Kalman gain, $\hat{x}_k^{-}$ and $P_k^{-}$ are
the prior (predicted) state and error variance, and $\hat{x}_k^{+}$ and
$P_k^{+}$ are the posterior (updated) estimates.

Additionally, the process noise $Q_m$ adapts dynamically to track rapid
changes: when the current innovation $|z_k - \hat{x}_k^{-}|$ exceeds 1.5
times the previous innovation, $Q_m$ is inflated by a factor of~1.8 (capped at
$0.5\,\sigma_m^2$), temporarily increasing the filter's responsiveness before
reverting to the baseline value.

\subsubsection*{Climatological prior database}

The adaptive Kalman filter described above relies on monthly climatological
statistics ($\mu_m$, $\sigma_m$, P01$_m$, P99$_m$) derived from a curated
database of 27 reference cities (Table~\ref{tab:ref_cities}). Each city's JSON
record stores monthly normals, standard deviations, and P01--P99 percentiles
for both ET\textsubscript{0} and precipitation, computed from the 1991--2020
period of the BR-DWGD dataset \citep{Xavier2016, Xavier2022} for Brazilian
sites, and from ERA5 reanalysis for the global sites.

\begin{table}[htbp]
\centering
\caption{Reference cities providing climatological priors for the adaptive
Kalman filter. Brazilian cities (region tag \texttt{brasil}) were also used
in the validation study; global cities provide priors for real-time queries
only and were not part of the validation.}
\label{tab:ref_cities}
\small
\begin{tabular}{@{}rlrrr@{}}
\toprule
\# & City & Lat.\ (\textdegree{}) & Lon.\ (\textdegree{}) & Elev.\ (m) \\
\midrule
\multicolumn{5}{l}{\textit{Brazil (17 cities --- MATOPIBA + Piracicaba)}} \\
1  & Alvorada do Gurgu\'eia, PI & $-8.44$ & $-43.86$ & 269 \\
2  & Aragua\'ina, TO            & $-7.19$ & $-48.21$ & 235 \\
3  & Balsas, MA                 & $-7.53$ & $-46.04$ & 256 \\
4  & Barreiras, BA              & $-12.14$ & $-45.00$ & 442 \\
5  & Bom Jesus, PI              & $-9.07$ & $-44.36$ & 300 \\
6  & Campos Lindos, TO          & $-7.97$ & $-46.80$ & 297 \\
7  & Carolina, MA               & $-7.33$ & $-47.47$ & 180 \\
8  & Corrente, PI               & $-10.44$ & $-45.16$ & 444 \\
9  & Formosa do Rio Preto, BA   & $-11.05$ & $-45.19$ & 496 \\
10 & Imperatriz, MA             & $-5.53$ & $-47.48$ & 131 \\
11 & Lu\'is Eduardo Magalh\~aes, BA & $-12.08$ & $-45.80$ & 769 \\
12 & Pedro Afonso, TO           & $-8.96$ & $-48.18$ & 182 \\
13 & Piracicaba, SP             & $-22.73$ & $-47.65$ & 524 \\
14 & Porto Nacional, TO         & $-10.71$ & $-48.41$ & 234 \\
15 & S\~ao Desid\'erio, BA      & $-12.36$ & $-44.97$ & 504 \\
16 & Tasso Fragoso, MA          & $-8.48$ & $-45.74$ & 240 \\
17 & Uru\c{c}u\'i, PI           & $-7.23$ & $-44.56$ & 175 \\
\midrule
\multicolumn{5}{l}{\textit{Global (10 cities --- real-time priors only)}} \\
18 & Addis Ababa, Ethiopia      &  $9.01$ &  $38.76$ & 2324 \\
19 & Des Moines (IA), USA       & $41.59$ & $-93.62$ & 294 \\
20 & Fresno (CA), USA           & $36.75$ & $-119.77$ & 100 \\
21 & Hanoi, Vietnam             & $21.03$ & $105.81$ & 16 \\
22 & Krasnodar, Russia          & $45.04$ &  $38.99$ & 34 \\
23 & Ludhiana, Punjab, India    & $30.90$ &  $75.86$ & 255 \\
24 & Mendoza, Argentina         & $-32.89$ & $-68.85$ & 827 \\
25 & Polokwane, Limpopo, S.\ Africa & $-23.90$ & $29.47$ & 1294 \\
26 & Seville, Spain             & $37.39$ &  $-5.98$ & 31 \\
27 & Wagga Wagga, Australia     & $-35.13$ & $147.35$ & 213 \\
\bottomrule
\end{tabular}
\end{table}

For a given query location, the system selects the nearest reference city
subject to three constraints enforced sequentially:

\begin{enumerate}
\item \textbf{Region segregation.} The query coordinates are classified into a
  geographic region (Brazil, USA, Europe, Africa, Asia, Oceania, or Global)
  using coordinate-based bounding boxes. Brazilian queries are matched
  exclusively against the 17 Brazilian reference cities, while non-Brazilian
  queries are matched exclusively against the 10 global cities.  This prevents
  climatological priors from an inappropriate climate zone (e.g.,\ tropical
  MATOPIBA) from being applied to a temperate query location, and vice versa.

\item \textbf{Hemisphere consistency.} The candidate city must lie in the same
  hemisphere as the query, with a $\pm5$\textdegree{} equatorial tolerance to
  accommodate locations near the equator.

\item \textbf{Distance threshold.} Among all candidates passing the above
  filters, the nearest city within 200~km (haversine distance) is selected.
\end{enumerate}

The outcome determines the Kalman filter mode:
\begin{itemize}
\item \textbf{High-precision mode} (reference found within 200~km): the
  \texttt{AdaptiveKalmanFilter} is instantiated with the full set of monthly
  priors ($\mu_m$, $\sigma_m$, P01$_m$, P99$_m$) from the matched city.  A
  separate filter instance is created per calendar month, ensuring
  season-specific initialisation.

\item \textbf{Global-fallback mode} (no reference within 200~km): a
  \texttt{SimpleKalmanFilter} with fixed parameters ($R = 0.8$, $Q = 0.05$,
  initial state $\hat{x}_0 = 5.0$~mm\,d$^{-1}$) provides temporal smoothing
  without climatological constraints.  This mode still improves on the raw
  Stage~1 output by suppressing day-to-day noise, but cannot perform the
  three-tier adaptive measurement noise inflation described in
  Eq.~\eqref{eq:adaptive_R}.
\end{itemize}


\subsection{FAO-56 Penman-Monteith ET\textsubscript{0} Calculation}
\label{subsec:eto}

Following data fusion, the daily reference evapotranspiration is computed using
the FAO-56 Penman-Monteith equation \citep{Allen1998}:

\begin{equation}
\text{ET}_0 = \frac{0.408\,\Delta\,(R_n - G) + \gamma\,\dfrac{900}{T_{\text{mean}}+273}\,u_2\,(e_s - e_a)}{\Delta + \gamma\,(1 + 0.34\,u_2)}
\label{eq:fao56}
\end{equation}

\noindent where $\text{ET}_0$ is the reference evapotranspiration
(mm\,d\textsuperscript{$-1$}), $R_n$ is the net radiation at the crop surface
(MJ\,m\textsuperscript{$-2$}\,d\textsuperscript{$-1$}), $G$ is the soil heat
flux density (MJ\,m\textsuperscript{$-2$}\,d\textsuperscript{$-1$}, assumed
zero for daily time steps), $T_{\text{mean}}$ is the mean daily air temperature
at 2~m height (\textdegree{}C), $u_2$ is the wind speed at 2~m height
(m\,s\textsuperscript{$-1$}), $e_s$ is the saturation vapour pressure (kPa),
$e_a$ is the actual vapour pressure (kPa), $\Delta$ is the slope of the
saturation vapour pressure curve (kPa\,\textdegree{}C\textsuperscript{$-1$}),
and $\gamma$ is the psychrometric constant
(kPa\,\textdegree{}C\textsuperscript{$-1$}).

The intermediate quantities are computed as specified in \citet{Allen1998}:

\textbf{Saturation vapour pressure:}
\begin{equation}
e_s = \frac{e^\circ(T_{\max}) + e^\circ(T_{\min})}{2}, \quad
e^\circ(T) = 0.6108\,\exp\!\left[\frac{17.27\,T}{T + 237.3}\right]
\label{eq:es}
\end{equation}

\textbf{Actual vapour pressure} (from mean relative humidity):
\begin{equation}
e_a = \frac{\text{RH}_{\text{mean}}}{100} \cdot e_s
\label{eq:ea}
\end{equation}

\textbf{Slope of the saturation vapour pressure curve:}
\begin{equation}
\Delta = \frac{4098 \cdot e^\circ(T_{\text{mean}})}{(T_{\text{mean}} + 237.3)^2}
\label{eq:slope}
\end{equation}

\textbf{Psychrometric constant:}
\begin{equation}
\gamma = 0.665 \times 10^{-3} \cdot P
\label{eq:gamma}
\end{equation}

\textbf{Net radiation:} The net shortwave radiation is $R_{ns} = (1 - \alpha)\,R_s$,
where $\alpha = 0.23$ is the grass reference surface albedo. The net longwave
radiation is computed using the Stefan-Boltzmann approach (FAO-56 Eq.~39):

\begin{equation}
R_{nl} = \sigma \left[\frac{T_{\max,K}^4 + T_{\min,K}^4}{2}\right]
\left(0.34 - 0.14\sqrt{e_a}\right)\,f
\label{eq:rnl}
\end{equation}

\noindent where $\sigma = 4.903 \times 10^{-9}$~MJ\,m$^{-2}$\,d$^{-1}$\,K$^{-4}$
is the Stefan-Boltzmann constant, $T_{\max,K}$ and $T_{\min,K}$ are the daily
maximum and minimum temperatures in Kelvin, and $f$ is the cloud cover factor:

\begin{equation}
f = 1.35\,\frac{R_s}{R_a} - 0.35
\label{eq:cloud_factor}
\end{equation}

\noindent bounded to the interval $[0.3,\, 1.0]$, where $R_a$ is the
extraterrestrial radiation. The net radiation is then:

\begin{equation}
R_n = R_{ns} - R_{nl} = (1 - \alpha)\,R_s - R_{nl}
\label{eq:rn}
\end{equation}

\textbf{Extraterrestrial radiation} ($R_a$) is computed from latitude and day
of year following the standard astronomical formulations in
\citet{Allen1998}, incorporating the inverse relative Earth--Sun distance,
solar declination ($\delta = 0.409\,\sin(2\pi\,J/365 - 1.39)$, where $J$ is
the day of year), and sunset hour angle with polar edge-case handling.

\textbf{Minimum physical guards.} Wind speed is clamped to $u_2 \geq 0.5$~m\,s$^{-1}$,
solar radiation to $R_s \geq 0.1$~MJ\,m$^{-2}$\,d$^{-1}$, and vapour pressure
deficit to $\text{VPD} \geq 0.01$~kPa to prevent numerical instabilities.

\textbf{Quality classification.} Each daily ET\textsubscript{0} estimate is
classified as ``high'' quality when $0.1 \leq \text{ET}_0 \leq 12$~mm\,d\textsuperscript{$-1$},
and ``low'' otherwise, enabling downstream filtering for anomalous values.

\subsubsection*{Water deficit analysis}

EVAonline computes the daily water deficit as:

\begin{equation}
\text{WD} = K_c \cdot \text{ET}_0 - P_{\text{eff}}
\label{eq:water_deficit}
\end{equation}

\noindent where $K_c$ is the crop coefficient (user-defined or default
$K_c = 1.0$ for reference grass) and $P_{\text{eff}}$ is the effective
precipitation derived from the fused precipitation data. The irrigation
requirement is $\text{IR} = \max(0, \text{WD})$, and the cumulative deficit
over the analysis period provides an estimate of the total irrigation volume
required.


\subsection{Software Design and Operational Modes}
\label{subsec:software}

\subsubsection*{Operational modes}

EVAonline supports three operational modes that determine which data sources
are queried and how the fusion pipeline is configured:

\textbf{(a) Recent mode (7 days).} Queries all available sources for the most
recent 7-day period. This mode provides the most robust fusion with maximum
source redundancy. NASA POWER data may exhibit a 2--3 day latency, compensated
by real-time sources (Open-Meteo Forecast, MET Norway, NWS).

\textbf{(b) Historical mode (custom date range).} Queries NASA POWER
(1981--present) and Open-Meteo Archive with ERA5 reanalysis (1940--present)
for user-specified historical periods. This mode enables long-term
climatological analyses and trend detection. Results are delivered
asynchronously via email as CSV or Excel file attachments, accommodating the
longer processing times required for multi-year queries.

\textbf{(c) Forecast mode (today + 5\,days = 6\,days).} Queries Open-Meteo
Forecast, MET Norway, and NWS Forecast (for US locations). This mode applies
the region-adaptive weights shown in Table~\ref{tab:weights}.

\subsubsection*{Asynchronous processing and real-time progress}

ET\textsubscript{0} calculations involve multiple sequential HTTP requests to
external climate APIs, data fusion computations, and database operations. To
prevent HTTP timeouts and provide a responsive user experience, all
calculations are processed asynchronously through the following pipeline:

\begin{enumerate}
\item The client submits a calculation request via HTTP POST.
\item The API creates an asynchronous Celery task and returns a unique task
  identifier (UUID) immediately.
\item The client establishes a WebSocket connection (via Redis pub/sub) to
  receive real-time progress updates.
\item The Celery worker executes the calculation pipeline, emitting progress
  events at each stage: elevation retrieval (10\%), climate data fetching
  (20--60\%, incremental per source), Kalman fusion (70\%), ET\textsubscript{0}
  computation (80\%), water deficit analysis (90\%), and completion (100\%).
\item The frontend renders a progress bar with step descriptions and displays
  the final results as interactive charts and tables.
\end{enumerate}

Failed tasks are retried up to three times with exponential backoff. The
system sends email notifications on both successful completion (with data file
attachment for historical mode) and on failure.

\subsubsection*{Caching strategy}

To minimize redundant API calls and reduce response times, EVAonline implements
an adaptive caching strategy using Redis with dynamic TTL based on data age
(Table~\ref{tab:cache}).

\begin{table}[htbp]
\centering
\caption{Redis caching strategy with dynamic time-to-live (TTL) based on the
temporal distance of the requested data from the current date.}
\label{tab:cache}
\begin{tabular}{@{}llrl@{}}
\toprule
Data category & Age criterion & TTL & Rationale \\
\midrule
Historical & $>$30 days old & 30 days & Stable reanalysis products \\
Recent & 7--30 days old & 1 day & Near-final reanalysis data \\
Very recent & $<$7 days old & 12 hours & May receive late corrections \\
Forecast & Future dates & 1 hour & NWP models update frequently \\
Elevation & Static & 24 hours & Terrain data, rarely changes \\
\bottomrule
\end{tabular}
\end{table}

\subsubsection*{User interface}

The interactive dashboard provides: (a)~a Leaflet.js map for click-to-select
coordinate input with geocoding, or manual latitude/longitude entry with
real-time validation; (b)~mode selection (Recent, Historical, Forecast) with
date range pickers; (c)~WebSocket-driven progress bars during calculation;
(d)~interactive Plotly charts for daily ET\textsubscript{0} time series,
monthly summaries, source-by-source comparison, and cumulative water deficit;
(e)~per-component CSV, Excel, and PNG download buttons; and (f)~complete
bilingual support with runtime English/Portuguese switching. A visitor tracking
system with session-based deduplication (Redis counters synchronized to
PostgreSQL) monitors platform usage.

\subsubsection*{REST API}

The API follows the OpenAPI~3.0 specification with automatic interactive
documentation (Swagger UI). Key endpoints include: \texttt{POST
/api/v1/eto/calculate} for submitting calculations, \texttt{GET
/api/v1/eto/result/\{task\_id\}} for retrieving results, \texttt{WS
/ws/eto/\{task\_id\}} for real-time progress, and \texttt{GET
/api/v1/health/detailed} for system health monitoring with individual service
status.


%%
%% RESULTS AND DISCUSSION
%%
\section{Results and Discussion}
\label{sec:results_discussion}

\subsection{Study Area and Validation Protocol}
\label{subsec:study_area}

The platform was validated using the MATOPIBA agricultural frontier region in
Brazil, one of the world's most dynamic cropland expansion areas. Seventeen
cities were selected across the four MATOPIBA states, spanning
latitudes from 5.52\textdegree{}S to 12.36\textdegree{}S and elevations from
100 to 900~m (Table~\ref{tab:sites}), ensuring representation of the region's
climatic and topographic variability. The city of Piracicaba (SP) was included
as an external reference site for cross-validation purposes, representing a
well-instrumented subtropical humid climate ($22.72\text{\textdegree{}S}$,
$47.63\text{\textdegree{}W}$).

\begin{table}[htbp]
\centering
\caption{Validation sites across the MATOPIBA region and reference site.
States: MA = Maranh\~ao, TO = Tocantins, PI = Piau\'i, BA = Bahia,
SP = S\~ao Paulo.}
\label{tab:sites}
\small
\begin{tabular}{@{}llcc@{}}
\toprule
City & State & Latitude (\textdegree{}S) & Longitude (\textdegree{}W) \\
\midrule
Alvorada do Gurgu\'eia & PI & 8.42 & 43.78 \\
Aragua\'ina & TO & 7.19 & 48.21 \\
Balsas & MA & 7.53 & 46.04 \\
Barreiras & BA & 12.15 & 45.00 \\
Bom Jesus & PI & 9.07 & 44.36 \\
Campos Lindos & TO & 7.99 & 46.87 \\
Carolina & MA & 7.33 & 47.47 \\
Corrente & PI & 10.44 & 45.16 \\
Formosa do Rio Preto & BA & 11.05 & 45.19 \\
Imperatriz & MA & 5.52 & 47.47 \\
Lu\'is Eduardo Magalh\~aes & BA & 12.10 & 45.80 \\
Pedro Afonso & TO & 8.97 & 48.17 \\
Porto Nacional & TO & 10.71 & 48.42 \\
S\~ao Desid\'erio & BA & 12.36 & 44.97 \\
Tasso Fragoso & MA & 8.47 & 45.75 \\
Uru\c{c}u\'i & PI & 7.23 & 44.56 \\
\midrule
Piracicaba (reference) & SP & 22.72 & 47.63 \\
\bottomrule
\end{tabular}
\end{table}

The validation protocol covers the 30-year period 1991--2020 (aligned with the
WMO standard climate normal period) and follows a seven-step procedure:
(1)~study area characterization with K\"oppen climate classification;
(2)~descriptive statistics of raw climate variables from each source;
(3)~data consolidation with standardized variable naming;
(4)~independent ET\textsubscript{0} calculation using NASA POWER only and
Open-Meteo Archive only;
(5)~cross-validation between individual-source ET\textsubscript{0} estimates;
(6)~execution of the complete two-stage fusion pipeline and comparison against
individual-source estimates; and
(7)~comprehensive multi-metric statistical comparison.

The reference dataset for validation is the Brazilian Daily Weather Gridded
Data (BR-DWGD) version~2 \citep{Xavier2016, Xavier2022}, which provides quality-controlled
daily meteorological data on a $0.1\text{\textdegree{}} \times 0.1\text{\textdegree{}}$
grid covering all of Brazil for the period 1961--2020.

Six statistical metrics are used (Table~\ref{tab:metrics_def}): Root Mean Square
Error (RMSE), Mean Absolute Error (MAE), Pearson correlation coefficient squared
($r^2$), Nash-Sutcliffe Efficiency \citep[NSE;][]{Nash1970}, Percent Bias (PBIAS),
and Kling-Gupta Efficiency \citep[KGE;][]{Gupta2009}.

\begin{table}[htbp]
\centering
\caption{Performance metrics used for validation.}
\label{tab:metrics_def}
\small
\begin{tabular}{@{}lll@{}}
\toprule
Metric & Formula & Optimal \\
\midrule
RMSE & $\sqrt{\frac{1}{n}\sum_{i=1}^{n}(E_i - O_i)^2}$ & 0 \\[6pt]
MAE & $\frac{1}{n}\sum_{i=1}^{n}|E_i - O_i|$ & 0 \\[6pt]
$r^2$ & $\left(\frac{\sum(O_i-\bar{O})(E_i-\bar{E})}{\sqrt{\sum(O_i-\bar{O})^2\sum(E_i-\bar{E})^2}}\right)^2$ & 1 \\[6pt]
NSE & $1 - \frac{\sum(E_i - O_i)^2}{\sum(O_i - \bar{O})^2}$ & 1 \\[6pt]
PBIAS (\%) & $\frac{\sum(E_i - O_i)}{\sum O_i} \times 100$ & 0 \\[6pt]
KGE & $1 - \sqrt{(r-1)^2 + (\beta-1)^2 + (\alpha-1)^2}$ & 1 \\
\bottomrule
\end{tabular}
\begin{flushleft}
\small $O_i$: observed; $E_i$: estimated; $r$: Pearson correlation;
$\beta$: bias ratio; $\alpha$: variability ratio.
\end{flushleft}
\end{table}


\subsection{Validation Results}
\label{subsec:results}

The two-stage fusion was validated against the BR-DWGD reference dataset
\citep{xavier2022} across 17 Brazilian sites (16 in the MATOPIBA region
plus Piracicaba/SP), totalling 186{,}286 daily observations over the period
1991--2020. Table~\ref{tab:results} summarises the overall performance of
each data source compared to the reference, while Figures~\ref{fig:scatter_nasa}--\ref{fig:seasonal} present detailed visual analyses.

\begin{table}[!ht]
\centering
\caption{Daily ET\textsubscript{0} validation results across 17 sites
(1991--2020, 186{,}286 daily observations). Values are per-site
means ($\pm$~standard deviation where shown). Pooled metrics across
the full dataset are reported in the scatter-plot captions
(Figures~\ref{fig:scatter_nasa}--\ref{fig:scatter_fusion}). Best values
in bold.}
\label{tab:results}
\small
\begin{tabular}{@{}lcccccc@{}}
\toprule
\textbf{Method} & \textbf{$R^2$} & \textbf{KGE} & \textbf{NSE} &
\textbf{MAE} & \textbf{RMSE} & \textbf{PBIAS} \\
 & & & & (mm\,d$^{-1}$) & (mm\,d$^{-1}$) & (\%) \\
\midrule
NASA POWER (FAO-56)
  & $0.740 \pm 0.062$ & $0.411 \pm 0.264$ & $-0.363 \pm 0.788$
  & $0.845$ & $1.117$ & $+15.78$ \\
Open-Meteo Archive (ERA5-Land)
  & $0.636 \pm 0.173$ & $0.432 \pm 0.413$ & $-0.547 \pm 1.820$
  & $0.859$ & $1.097$ & $+13.02$ \\
Open-Meteo API (ERA5)
  & $0.649 \pm 0.174$ & $0.584 \pm 0.188$ & $0.216 \pm 0.356$
  & $0.690$ & $0.860$ & $+8.27$ \\
\textbf{EVAonline Fusion}
  & $\mathbf{0.694 \pm 0.074}$ & $\mathbf{0.814 \pm 0.053}$
  & $\mathbf{0.676 \pm 0.085}$ & $\mathbf{0.423}$
  & $\mathbf{0.566}$ & $\mathbf{+0.71}$ \\
\midrule
\multicolumn{7}{l}{\footnotesize Improvement vs.\ best individual source (Open-Meteo API, KGE\,=\,0.584):} \\
\multicolumn{7}{l}{\footnotesize $\Delta$KGE: +39\% $|$ $\Delta$MAE: $-$39\% $|$ $\Delta$RMSE: $-$34\% $|$ $\Delta$PBIAS: $-$91.5\%} \\
\multicolumn{7}{l}{\footnotesize Improvement vs.\ NASA POWER (most common single source):} \\
\multicolumn{7}{l}{\footnotesize $\Delta$KGE: +98\% $|$ $\Delta$MAE: $-$50\% $|$ $\Delta$RMSE: $-$49\% $|$ $\Delta$PBIAS: $-$95.5\%} \\
\bottomrule
\end{tabular}
\end{table}

% ======================= FIGURE (a): NASA POWER SCATTER =======================
\begin{figure}[!ht]
\centering
\includegraphics[width=0.85\textwidth]{Figure_6a_scatter_nasa.png}
\caption{Density scatter plot of NASA POWER ET\textsubscript{0}
vs.\ BR-DWGD reference (186{,}286 daily observations, 17 sites,
1991--2020). Red dashed line: 1:1 reference; blue solid line: OLS
regression ($y = 1.343x - 0.785$; pooled $R^2 = 0.737$,
KGE\,=\,0.396).}
\label{fig:scatter_nasa}
\end{figure}

% ==================== FIGURE (b): OPEN-METEO ARCHIVE SCATTER ==================
\begin{figure}[!ht]
\centering
\includegraphics[width=0.85\textwidth]{Figure_6b_scatter_openmeteo.png}
\caption{Density scatter plot of Open-Meteo Archive (ERA5-Land)
ET\textsubscript{0} vs.\ BR-DWGD reference. OLS regression:
$y = 1.096x + 0.137$; pooled $R^2 = 0.538$, KGE\,=\,0.423,
RMSE\,=\,1.229~mm\,d$^{-1}$.}
\label{fig:scatter_openmeteo}
\end{figure}

% ==================== FIGURE (c): EVAONLINE FUSION SCATTER ====================
\begin{figure}[!ht]
\centering
\includegraphics[width=0.85\textwidth]{Figure_6c_scatter_fusion.png}
\caption{Density scatter plot of EVAonline two-stage fusion
ET\textsubscript{0} vs.\ BR-DWGD reference. OLS regression:
$y = 0.819x + 0.807$; pooled KGE\,=\,0.849, NSE\,=\,0.720,
RMSE\,=\,0.570~mm\,d$^{-1}$, PBIAS\,=\,+0.70\%.}
\label{fig:scatter_fusion}
\end{figure}

% =================== FIGURE (d): ERROR DISTRIBUTION BOXPLOTS ==================
\begin{figure}[!ht]
\centering
\includegraphics[width=0.85\textwidth]{Figure_6d_error_distribution.png}
\caption{Box plots of daily estimation error
(estimated~$-$~reference) for NASA POWER, Open-Meteo Archive, and
EVAonline Fusion. Median errors: $+0.69$, $+0.55$, and
$+0.03$~mm\,d$^{-1}$, respectively. Fliers are omitted for clarity.}
\label{fig:error_dist}
\end{figure}

% ====================== FIGURE (e): KGE BY CITY ==============================
\begin{figure}[!ht]
\centering
\includegraphics[width=0.95\textwidth]{Figure_6e_kge_by_city.png}
\caption{Kling-Gupta Efficiency (KGE) by site for all four data
sources. The EVAonline Fusion (green) achieves KGE~$> 0.72$ at every
site, whereas individual sources show large inter-site variability and
negative values at several locations.}
\label{fig:kge_city}
\end{figure}

% ==================== FIGURE (f): SEASONAL PERFORMANCE ========================
\begin{figure}[!ht]
\centering
\includegraphics[width=0.90\textwidth]{Figure_6f_seasonal_performance.png}
\caption{Monthly RMSE for NASA POWER, Open-Meteo Archive, and
EVAonline Fusion. The shaded band marks the dry season
(May--September).}
\label{fig:seasonal}
\end{figure}

Several key findings emerge from the validation analysis.

\textbf{Individual-source characterisation.}
The pooled scatter plots (Figures~\ref{fig:scatter_nasa}
and~\ref{fig:scatter_openmeteo}) reveal distinct error structures for
each source. NASA POWER shows a regression slope well above unity
($y = 1.343x - 0.785$), indicating overestimation that intensifies at
higher ET\textsubscript{0} values; despite a pooled $R^2 = 0.737$, the
large bias ($+0.69$~mm\,d$^{-1}$, PBIAS\,=\,+16.1\%) depresses the
KGE to only 0.396. Open-Meteo Archive displays greater scatter
($R^2 = 0.538$) but a nearly unit regression slope
($y = 1.096x + 0.137$), suggesting a dominant constant-offset
overestimation component (PBIAS\,=\,+12.8\%); the pooled RMSE reaches
1.229~mm\,d$^{-1}$ and KGE is 0.423. Although individually limited,
the two sources exhibit complementary error structures---a property
exploited by the fusion framework.

\textbf{Systematic bias elimination.}
The most striking result of the
two-stage fusion is the near-complete elimination of systematic error.
Individual sources exhibited substantial positive biases: $+15.78\%$
(NASA POWER), $+13.02\%$ (Open-Meteo Archive), and $+8.27\%$ (Open-Meteo
API). The fusion reduced PBIAS to only $+0.71 \pm 0.53\%$---a 95.5\%
reduction relative to NASA POWER and 91.5\% relative to Open-Meteo API
(Figure~\ref{fig:scatter_fusion}). The error-distribution box
plots (Figure~\ref{fig:error_dist}) confirm this: NASA POWER and
Open-Meteo Archive present median errors of $+0.69$ and
$+0.55$~mm\,d$^{-1}$, respectively, with wide interquartile ranges,
whereas the fusion compresses the median error to
$+0.03$~mm\,d$^{-1}$ and reduces the interquartile range by more than
50\%. This near-zero bias is critical for cumulative water-balance
applications, where systematic errors compound over time, and for
irrigation scheduling, where persistent overestimation of
ET\textsubscript{0} would lead to excessive water allocation.

\textbf{Error reduction.}
The pooled fusion scatter plot (Figure~\ref{fig:scatter_fusion})
shows markedly tighter clustering around the 1:1 line, with a pooled
KGE of 0.849, NSE of 0.720, RMSE of 0.570~mm\,d$^{-1}$, and near-zero
bias ($+0.030$~mm\,d$^{-1}$, PBIAS\,=\,+0.70\%). The regression line
($y = 0.819x + 0.807$) indicates slight underestimation at higher
ET\textsubscript{0} values, balanced by the intercept offset. In
per-site terms, the fusion reduced RMSE by 49.3\% compared to
NASA POWER alone (from 1.117 to 0.566~mm\,d$^{-1}$) and by 48.3\%
compared to Open-Meteo Archive (from 1.097 to 0.566~mm\,d$^{-1}$).
Relative to the best individual source (Open-Meteo API, RMSE\,=\,0.860),
the reduction was 34.2\%. MAE was reduced by 50\% compared to NASA POWER
(from 0.845 to 0.423~mm\,d$^{-1}$) and by 39\% compared to Open-Meteo
API (from 0.690 to 0.423~mm\,d$^{-1}$).

\textbf{Seasonal complementarity.}
Performance improvement was most
pronounced during the dry season (May--September), when individual sources
exhibited highest errors (Figure~\ref{fig:seasonal}). During the dry months,
fusion RMSE was 0.49~mm\,d$^{-1}$ versus 1.46~mm\,d$^{-1}$ for NASA POWER
(66\% reduction) and 1.50~mm\,d$^{-1}$ for Open-Meteo (67\% reduction).
During the wet season, the fusion maintained RMSE of 0.62~mm\,d$^{-1}$
versus 0.90~mm\,d$^{-1}$ (NASA) and 0.99~mm\,d$^{-1}$ (Open-Meteo),
representing reductions of 31\% and 37\%, respectively. The adaptive Kalman
filter effectively exploited inter-source complementarity by adjusting
its innovation weighting to the climatological context of each month.

\textbf{Spatial consistency.}
The fusion achieved KGE~$>0.72$ at all
17 validation sites (Figure~\ref{fig:kge_city}), with a coefficient of
variation of only 6.5\% ($0.814 \pm 0.053$), compared to 32\% for
Open-Meteo API ($0.584 \pm 0.188$) and 64\% for NASA POWER
($0.411 \pm 0.264$). Notably, several individual-source estimates
exhibit negative KGE at specific sites (e.g.,\ Open-Meteo Archive at
Campos Lindos, KGE\,=\,$-0.61$), meaning that a constant climatological
mean would outperform those single-source estimates; the fusion recovers
acceptable performance even at these difficult locations. This spatial
consistency is particularly important for regional-scale applications in
the MATOPIBA agricultural frontier, where station networks are sparse
and users require reliable estimates regardless of specific location.

\textbf{Radiation dominance.} Solar radiation ($R_s$) emerged as the most
critical variable for ET\textsubscript{0} accuracy. For historical periods,
NASA POWER's satellite-derived CERES radiation product was typically the only
available source with daily resolution, entering the fusion at 100\% weight
(single-source mode). For forecast and recent modes, multi-source
fusion of $R_s$ provided additional robustness.

\textbf{Climatological constraint benefit.} Sites within the 200~km
radius of reference cities benefited from the full adaptive Kalman filter
with P01/P99 climatological bounds. The adaptive measurement noise
inflation effectively rejected anomalous values that would have degraded
daily estimates. Sites falling back to the simplified (fixed-parameter)
Kalman filter showed less improvement, underscoring the value of regional
climatological priors.


\subsection{Comparison with Existing Tools}
\label{subsec:comparison}

Table~\ref{tab:comparison} compares EVAonline with established
ET\textsubscript{0} estimation tools across key functional dimensions.

\begin{table}[htbp]
\centering
\caption{Feature comparison of EVAonline with existing ET\textsubscript{0}
estimation tools.}
\label{tab:comparison}
\small
\begin{tabular}{@{}lcccccc@{}}
\toprule
Feature & \rotatebox{60}{EVAonline} & \rotatebox{60}{CROPWAT 8.0} & \rotatebox{60}{ETo Calculator} & \rotatebox{60}{PyETo} & \rotatebox{60}{Ref-ET} & \rotatebox{60}{OpenET} \\
\midrule
Web-based (no install) & \checkmark & --- & --- & --- & --- & \checkmark \\
Multi-source data fusion & \checkmark & --- & --- & --- & --- & \checkmark$^a$ \\
Automatic data acquisition & \checkmark & --- & --- & --- & --- & \checkmark \\
FAO-56 Penman-Monteith & \checkmark & \checkmark & \checkmark & \checkmark & \checkmark & ---$^b$ \\
Global coverage & \checkmark & manual & manual & manual & manual & USA only \\
Open source & AGPL-3.0 & --- & --- & BSD & --- & partial \\
Real-time visualization & \checkmark & --- & --- & --- & --- & \checkmark \\
Water deficit analysis & \checkmark & \checkmark & --- & --- & --- & --- \\
API access & REST & --- & --- & Python & --- & REST \\
Free of charge & \checkmark & \checkmark & \checkmark & \checkmark & --- & \checkmark \\
Docker deployment & \checkmark & --- & --- & pip & --- & --- \\
Bilingual interface & EN/PT & multi & EN & EN & EN & EN \\
\bottomrule
\end{tabular}
\begin{flushleft}
\small $^a$OpenET fuses satellite-based actual ET models, not meteorological
data.\\$^b$OpenET estimates actual ET from satellite imagery, not
reference ET\textsubscript{0}.
\end{flushleft}
\end{table}

EVAonline is, to the authors' knowledge, the only freely available tool that
combines automatic multi-source meteorological data acquisition, optimal data
fusion with quantified uncertainty through Kalman filtering, real-time web-based
access with interactive visualization, and global coverage in a single
integrated platform. CROPWAT and the ETo Calculator remain important tools for
desktop use but require manual data entry, offering no data fusion capability.
PyETo provides a clean programmatic interface but targets developers rather than
end users and offers no data acquisition or visualization. OpenET
provides satellite-based actual ET for the United States but does not estimate
reference ET\textsubscript{0} and is limited in geographic coverage.


\subsection{Strengths, Limitations, and Future Directions}
\label{subsec:limitations}

\subsubsection*{Strengths}

The main strengths of EVAonline include: (a)~the two-stage fusion strategy
that exploits the complementary strengths of multiple data sources while
mitigating their individual weaknesses; (b)~the adaptive Kalman filter with
climatological priors that provides physically informed quality control beyond
simple averaging; (c)~the hexagonal architecture that enables straightforward
addition of new data sources through the standardized adapter interface;
(d)~the Docker-based deployment that eliminates installation complexity and
ensures reproducibility across environments; and (e)~the fully open-source
AGPL-3.0 license that enables community contributions, peer review, and
institutional deployment.

\subsubsection*{Limitations}

Several limitations should be acknowledged. First, the cloud cover factor
in the Stefan-Boltzmann net longwave calculation uses $R_s/R_a$ rather
than $R_s/R_{so}$ (where $R_{so} = (0.75 + 2 \times 10^{-5}\,z)\,R_a$
includes the elevation correction); for low-to-moderate elevations
($z < 1000$~m) this introduces negligible error, but at high altitudes
a small overestimation of $R_{nl}$ and corresponding underestimation of
$R_n$ may occur. Second, the
climatological priors for the adaptive Kalman filter are currently available for
27 reference cities; locations beyond 200~km from any reference site fall back
to a simplified filter without climatological constraints, resulting in reduced
filtering effectiveness. Third, the platform currently relies exclusively on
gridded and reanalysis data products; integration of ground-based weather
station networks (e.g., INMET for Brazil, GHCN-Daily globally) would provide
additional independent observations to improve fusion quality. Fourth, the
ET\textsubscript{0} calculation is implemented as a scalar per-record
computation using Python's \texttt{math} module; while the data preprocessing
pipeline and Kalman filters leverage NumPy for array-level operations, fully
vectorizing the FAO-56 PM calculation would improve computational throughput
for large historical queries.

\subsubsection*{Future directions}

Priority areas for future development include: (a)~integration of ground-based
observation networks as additional fusion sources; (b)~expansion of the
climatological prior database from 27 to 200+ reference locations to achieve
near-global coverage with adaptive Kalman filtering; (c)~implementation of an
Extended Kalman Filter (EKF) that operates directly on the nonlinear FAO-56 PM
equation, fusing meteorological observations and ET\textsubscript{0} estimation
in a single step; (d)~machine learning-based bias correction layers trained on
regional station data; and (e)~development of ensemble-based uncertainty
quantification that provides confidence intervals for the final
ET\textsubscript{0} estimates, enabling risk-aware irrigation scheduling.


%%
%% CONCLUSIONS
%%
\section{Conclusions}
\label{sec:conclusions}

This paper presented EVAonline, an open-source, web-based platform for
reference evapotranspiration estimation that addresses the persistent challenge
of meteorological data scarcity and source-specific biases through multi-source
data fusion. The main contributions and findings are:

\begin{enumerate}
\item \textbf{Two-stage fusion methodology.} The combination of region-adaptive
weighted averaging (Stage~1) with an adaptive Kalman filter constrained by
30-year climatological priors (Stage~2) achieved a KGE of $0.814 \pm 0.053$
when validated against the BR-DWGD reference dataset at 17 MATOPIBA sites.
The fusion reduced RMSE by 34\% compared to the best individual source and
produced a near-zero bias (PBIAS~$= 0.71\%$), demonstrating that systematic
integration of complementary data sources yields measurably superior
ET\textsubscript{0} estimates.

\item \textbf{Accessibility.} By packaging the entire platform as
Docker-containerized microservices accessible through a web browser, EVAonline
eliminates technical barriers that have historically limited access to
calibrated ET\textsubscript{0} estimation. No local software installation,
programming knowledge, API keys, or data preprocessing are required.

\item \textbf{Extensibility.} The hexagonal architecture and standardized
adapter interface make it straightforward to integrate additional climate data
sources as they become available. The Kalman filter framework naturally
accommodates sources with different spatial resolutions, temporal coverages,
and quality levels.

\item \textbf{Reproducibility.} The complete source code (AGPL-3.0), Docker
deployment configuration, and independent validation dataset (published on
Zenodo) enable full reproduction of the results presented in this paper.
\end{enumerate}

EVAonline demonstrates that Kalman-filter-based data fusion of freely available
global climate data can provide reliable ET\textsubscript{0} estimates for
data-scarce agricultural regions, contributing to improved irrigation management
and water resource planning. The platform is particularly relevant for regions
such as MATOPIBA where agricultural expansion is outpacing the development of
meteorological observation infrastructure.


%%
%% Supporting sections
%%

\section*{CRediT authorship contribution statement}
\label{sec:credit_authorship}

%% TODO: Fill in based on actual contributions. Example format:
%% \textbf{Angela Cunha Soares}: Conceptualization, Methodology, Software,
%% Validation, Data Curation, Writing -- Original Draft, Writing -- Review
%% \& Editing, Visualization, Project Administration.
%% \textbf{Second Author}: Supervision, Writing -- Review \& Editing.


\section*{Software and data}
\label{sec:software_data}

The EVAonline source code is available under the AGPL-3.0-or-later license at
\url{https://github.com/angelacunhasoares/EVAONLINE}. The complete validation
dataset, analysis scripts, and reproducible Jupyter notebooks are published as
an independent, citable package on Zenodo:

%% TODO: Insert Zenodo DOI when minted
%% \url{https://doi.org/10.5281/zenodo.XXXXXXX}

The validation package includes its own \texttt{environment.yml} (Conda) and
\texttt{requirements.txt} (pip) for reproducibility independent of the main
application.


\section*{Declaration of competing interest}
\label{sec:declaration_competing_interest}

The authors declare no conflict of interest relevant to this work.


\section*{Acknowledgments}
\label{sec:acknowledgments}

%% TODO: Add funding agencies, computing resources, and any other
%% acknowledgments as appropriate. Example:
%% This work was supported by [funding agency] under grant [number].
%% The authors thank the NASA POWER, Open-Meteo, MET Norway, and NWS teams
%% for providing free access to their climate data APIs.


%%
%% APPENDIX
%%
\appendix

\section{Physical Limits for Climate Data Quality Control}
\label{app:limits}

Table~\ref{tab:limits_global} presents the global and Brazil-specific physical
limits used in the quality control stage. These bounds are derived from official
WMO records, NOAA documentation, and the BR-DWGD dataset
\citep{Xavier2016, Xavier2022}.

\begin{table}[htbp]
\centering
\caption{Physical range bounds for climate data quality control. Values outside
these ranges are set to NaN before entering the fusion pipeline.}
\label{tab:limits_global}
\small
\begin{tabular}{@{}lcccc@{}}
\toprule
\multirow{2}{*}{Variable} & \multicolumn{2}{c}{Global limits} & \multicolumn{2}{c}{Brazil limits} \\
\cmidrule(lr){2-3}\cmidrule(l){4-5}
 & Min & Max & Min & Max \\
\midrule
Temperature (\textdegree{}C) & $-90$ & 60 & $-30$ & 50 \\
Relative humidity (\%) & 0 & 100 & 0 & 100 \\
Wind speed (m\,s$^{-1}$) & 0 & 113 & 0 & 100 \\
Solar radiation (MJ\,m$^{-2}$\,d$^{-1}$) & 0 & 45 & 0 & 40 \\
Precipitation (mm\,d$^{-1}$) & 0 & 2000 & 0 & 450 \\
Atmospheric pressure (hPa) & 800 & 1150 & 900 & 1100 \\
\bottomrule
\end{tabular}
\begin{flushleft}
\small Global records: Death Valley 56.7~\textdegree{}C (2021), Vostok
$-89.2$~\textdegree{}C (1983), Bridge Creek tornado 113.3~m\,s$^{-1}$
(1999), Cilaos 1825~mm\,d$^{-1}$ (1952).
\end{flushleft}
\end{table}


\section{REST API Endpoints}
\label{app:api}

Table~\ref{tab:api} lists the principal REST API endpoints exposed by the
EVAonline backend.

\begin{table}[htbp]
\centering
\caption{Principal EVAonline REST API endpoints.}
\label{tab:api}
\small
\begin{tabular}{@{}llp{6.5cm}@{}}
\toprule
Method & Endpoint & Description \\
\midrule
POST & \texttt{/api/v1/eto/calculate} & Submit an ET\textsubscript{0} calculation request \\
GET & \texttt{/api/v1/eto/result/\{task\_id\}} & Retrieve calculation results \\
WS & \texttt{/ws/eto/\{task\_id\}} & Real-time progress stream via WebSocket \\
GET & \texttt{/api/v1/health/detailed} & System health and service status \\
GET & \texttt{/api/v1/climate/sources/available} & List available climate data sources \\
GET & \texttt{/api/v1/visitors/count} & Platform usage statistics \\
\bottomrule
\end{tabular}
\begin{flushleft}
\small Full API documentation is available via Swagger UI at
\texttt{/api/v1/docs}.
\end{flushleft}
\end{table}


%%
%% REFERENCES
%%
\bibliographystyle{elsarticle-harv}
\bibliography{EMS}

\end{document}
